\SeventhArticle{13}
To determine the degree of accuracy of our formula for the integral, $RA+R'A'+R''A''+\text{etc.}+R^{(n)}A^{(n)}$, let us set in general
\[
Ra^m+R'a'{}^m+R''a''{}^m+\text{etc.}+R^{(n)}a^{(n)m}=\frac{1}{m+1}-k^{(m)}
\]
so that $k^{(m)}$ is the difference between the true and the approximate value of the integral $\int t^m\,\mathrm{d}t$ taken from $t=0$ to $t=1$. Thus, on expanding each fraction into a series, we shall have\SourceLeaf{190}
\[
\begin{aligned}
&\frac{R}{t-a}+\frac{R'}{t-a'}+\frac{R''}{t-a''}+\text{etc.}+\frac{R^{(n)}}{t-a^{(n)}}\\
&\quad=(1-k)t^{-1}+(\tfrac12-k')t^{-2}+(\tfrac13-k'')t^{-3}\\ &\qquad+(\tfrac14-k''')t^{-4}+\text{etc.}\\
&\quad=t^{-1}+\tfrac12t^{-2}+\tfrac13t^{-3}+\tfrac14t^{-4}+\text{etc.}-\theta
\end{aligned}
\]
if we set
\[
\theta=kt^{-1}+k't^{-2}+k''t^{-3}+k'''t^{-4}+\text{etc.}
\]
or rather (since we already know that $k$, $k'$, $k''$, $k'''$, etc., up to $k^{(n)}$, must vanish automatically)
\[
\theta=k^{(n+1)}t^{-(n+2)}+k^{(n+2)}t^{-(n+3)}+k^{(n+3)}t^{-(n+4)}+\text{etc.}
\]

Multiplication by $T$ gives
\[
\begin{aligned}&T\left(\frac{R}{t-a}+\frac{R'}{t-a'}+\frac{R''}{t-a''}+\text{etc.}+\frac{R^{(n)}}{t-a^{(n)}}\right)\\ &=T'+T''-T\theta\end{aligned}
\]
The left side of this equation is a polynomial in $t$ of degree $n$, and its values for $t=a$, $t=a'$, $t=a''$, etc., are respectively $MR$, $M'R'$, $M''R''$, etc. The same holds for the function $T'$, as is evident from the very way in which the numbers $R$, $R'$, $R''$, etc., are determined. Thus the left side of the equation must be identical with $T'$, and hence $T''=T\theta$. Therefore $\theta$ is obtained by expanding the fraction $\frac{T''}{T}$; in this way the coefficients $k^{(n+1)}$, $k^{(n+2)}$, etc., can be determined as far as desired. Once these are found, the correction to our approximate value of the integral $\int y\,\mathrm{d}t$ will be
\[
=k^{(n+1)}K^{(n+1)}+k^{(n+2)}K^{(n+2)}+\text{etc.}
\]
if the series into which $y$ is expanded is
\[
y=K+K't+K''tt+K'''t^3+\text{etc.}
\]

\SeventhArticle{14}
If one prefers to express the correction in terms of the coefficients of the series in powers of $t-\tfrac12$,
\[
y=L+L'(t-\tfrac12)+L''(t-\tfrac12)^2+L'''(t-\tfrac12)^3+\text{etc.}
\]
it will be\SourceLeaf{191}
\[
=l^{(n+1)}L^{(n+1)}+l^{(n+2)}L^{(n+2)}+l^{(n+3)}L^{(n+3)}+\text{etc.}
\]
if, in general, $l^{(m)}$ denotes the correction to the approximate value of the integral $\int(t-\tfrac12)^m\,\mathrm{d}t$. These corrections $l^{(m)}$ will be connected with the corrections $k^{(m)}$ by the equation
\[
\begin{aligned}&l^{(m)}=k^{(m)}-\tfrac12mk^{(m-1)}+\tfrac14\cdot\frac{m\cdot m-1}{1\cdot2}k^{(m-2)}\\ &-\tfrac18\cdot\frac{m\cdot m-1\cdot m-2}{1\cdot2\cdot3}k^{(m-3)}+\text{etc.}\end{aligned}
\]
To find them independently, however, let us consider that the function $\theta$, under the substitution $t=\tfrac12u+\tfrac12$, becomes
\[
\begin{aligned}
&\phantom{{}+{}}2k(u^{-1}-u^{-2}+u^{-3}-u^{-4}+\text{etc.})\\
&+4k'(u^{-2}-2u^{-3}+3u^{-4}-4u^{-5}+\text{etc.})\\
&+8k''(u^{-3}-3u^{-4}+6u^{-5}-10u^{-6}+\text{etc.})\\
&+16k'''(u^{-4}-4u^{-5}+10u^{-6}-20u^{-7}+\text{etc.})\\
&+\text{etc.}
\end{aligned}
\]
or
\[
\begin{aligned}
&2ku^{-1}+4(k'-\tfrac12)u^{-2}+8(k''-\tfrac12\cdot2k'+\tfrac14k)u^{-3}\\
&\qquad+16(k'''-\tfrac12\cdot3k''+\tfrac14\cdot3k'-\tfrac18k)u^{-4}+\text{etc.}
\end{aligned}
\]
or
\[
2lu^{-1}+4l'u^{-2}+8l''u^{-3}+16l'''u^{-4}+\text{etc.}
\]
or finally, since we know in advance that $l$, $l'$, $l''$, $l'''$, etc., up to $l^{(n)}$, vanish automatically,
\[
\begin{aligned}&2^{n+2}l^{(n+1)}u^{-(n+2)}+2^{n+3}l^{(n+2)}u^{-(n+3)}\\ &+2^{n+4}l^{(n+4)}u^{-(n+4)}+\text{etc.}\end{aligned}
\]
But $\theta=\frac{T''}{T}$. Since $T$, $T''$, under the substitution $t=\tfrac12u+\tfrac12$, become $\frac{U}{2^{n+1}}$, $\frac{U''}{2^n}$ (Article 9), the function $\theta$ becomes $\frac{2U''}{U}$ under the same substitution. Thus, if we denote the series obtained by expanding the fraction $\frac{U''}{U}$ by $\Omega$, we shall have
\[
\begin{aligned}&\Omega=2^{n+1}l^{(n+1)}u^{-(n+2)}+2^{n+2}l^{(n+2)}u^{-(n+3)}\\ &+2^{n+3}l^{(n+3)}u^{-(n+4)}+\text{etc.}\end{aligned}
\]
In this way the coefficients $l^{(n+1)}$, $l^{(n+2)}$, etc., can be found as far as desired.

Thus, in the example of Article 10, we find
\[
\begin{aligned}
U''&=-\tfrac{176}{13125}u^{-1}-\tfrac{304}{28125}u^{-3}-\tfrac{2576}{309375}u^{-5}-\text{etc.}\\
\Omega&=-\tfrac{176}{13125}u^{-7}-\tfrac{832}{28125}u^{-9}-\tfrac{189856}{4296875}u^{-11}-\text{etc.}
\end{aligned}
\]
\SourceLeaf{192}and therefore the correction to the approximate value of the integral is
\[
=-\frac{11}{52500}L^{\mathrm{VI}}-\frac{13}{112500}L^{\mathrm{VIII}}-\frac{5933}{137500000}L^{\mathrm{X}}-\text{etc.}
\]
\SeventhArticle{15}
The coefficient $K^{(m)}$ in the series expansion of the function $y$ is, by \mbox{\textsc{Taylor}}’s theorem, equal to the value of
\[
\frac{1}{1.2.3\mathinner{.\,.\,.\,.}m}\mathbin{.}\frac{\mathrm{d}^m y}{\mathrm{d}t^m}
\quad\text{or}\quad
\frac{\Delta^m}{1.2.3\mathinner{.\,.\,.\,.}m}\mathbin{.}\frac{\mathrm{d}^m y}{\mathrm{d}x^m}
\]
for $t=0$, or $x=g$. Likewise, the coefficient $L^{(m)}$ is the value of the same expression for $t=\frac12$, or $u=0$, or $x=g+\frac12\Delta$. We shall assign \textit{order} $m$ to each coefficient. Generally speaking, therefore, our integration will be exact through order $n$, inclusive, whatever values are taken for $a,a',a''\mathinner{.\,.\,.\,.}a^{(n)}$. This, however, does not prevent a higher degree of accuracy from being attained by a suitable choice of these values. Thus we have already seen above that in the method of \textsc{Cotes}, that is, for $a=0,a'=\frac1n,a''=\frac2n,a'''=\frac3n$, etc., the accuracy automatically extends through order $n+1$, inclusive, whenever $n$ is an even number. In general, it is clear that if the values $a,a',a'',a'''$, etc., are chosen so that one or more initial terms of the function $T''$ or $U''$ vanish, the accuracy will advance beyond order $n$ by as many orders as there are terms that vanish. Hence it is easily inferred that, since the number of quantities we are free to choose is $n+1$, their suitable determination can always raise the accuracy through order $2n+1$, inclusive. In this way, with the aid of $n+1$ terms, it will be possible to attain the same order of accuracy that would require $2n+1$ or $2n+2$ terms if we followed the method of \textsc{Cotes}.

